\documentclass[10pt]{article}
\usepackage[margin=0.85in]{geometry}
\usepackage{amsmath,booktabs,hyperref,xurl}
\usepackage[T1]{fontenc}
\setlength{\parindent}{0pt}
\setlength{\parskip}{5pt}
\title{Why the 20-cm-margin frame fails model agreement}
\author{Collision Avoidance Research}
\date{October 3, 2026}
\begin{document}
\maketitle
\textbf{Finding.} The initialization is dynamically consistent, but is not
sufficiently compatible with the tightened constraints to support the current
single-linearization correction and damping rule. The fresh potential-field
seed misses the new sampled clearance requirement by 2.464 cm. Restoring its
convex feasibility already produces more nonlinear prediction error than the
10-mm acceptance threshold. This is a local numerical-design limitation, not
proof that avoiding the target is physically impossible.

\textbf{Scope and reproduction.} The 15-m/s accelerating head-on case stops
at $t=0.25$ s before its sixth control. The frozen continuation from
\path{/home/zai/.cache/collisionAvoidance/margin-20cm-20261003/validated-campaign/speed15-acceleratingHeadOn.mat}
was replayed through a temporary diagnostic copy of
\path{controller/solvePredictiveControl.m}. Only the copied function name,
output-directory argument and passive candidate saves changed. The online
controller and all parameters remain unchanged. The captures reproduce both
original disagreement values to $10^{-10}$ m. Results are offline numerical
model comparisons, not an additional closed-loop experiment or timing test.

\textbf{What the error measures.} For the optimized affine state and a nonlinear
RK4 rollout of exactly the same input sequence, the controller evaluates
\[
 E_i=\|p_i^{\rm aff}-p_i^{\rm nl}\|_2+
 R\,|\psi_i^{\rm aff}-\psi_i^{\rm nl}|,
 \qquad R=\sqrt{(4.8/2)^2+(1.9/2)^2}.
\]
The maximum over collision-relevant prediction nodes must not exceed 0.01 m;
state and CLF errors have separate gates. This is not the discrepancy at the
currently measured state. Both anchors are nonlinear bicycle rollouts and have
exactly zero discrepancy before input correction.

\begin{center}
\begin{tabular}{lrrr}
\toprule
Point & Node gap (cm) & First-step error (mm) & Maximum error (mm) \\
\midrule
Shifted anchor & 19.8814 & 0 & 0 \\
Shifted CLF solution & 19.8841 & 0.107674 & 10.534877 \\
Fresh potential-field anchor & 17.5362 & 0 & 0 \\
Fresh PCBF solution & 19.8340 & 0.066526 & 17.439265 \\
Fresh CLF solution & 19.7737 & 0.000970 & 26.946763 \\
\bottomrule
\end{tabular}
\end{center}
Node gaps above use nonlinear RK4 states at starts and their interpolated
midpoints, matching the controller's numerical geometry. They are not refined
continuous-time clearances. The 0.20-m inequalities hold for the optimized
\emph{affine} models; nonlinear rollout can have a smaller gap.

\textbf{Shifted trajectory.} The old inputs, rerolled from the current measured
state, miss the 20-cm collision row by 1.186 mm and also violate the terminal
membership cone by 0.041818 in that cone's units (not meters). Therefore the
zero-correction shifted anchor is not a feasible base for convex interpolation.
The CLF solution is convex-feasible but has 10.535-mm maximum disagreement,
only 0.535 mm beyond the gate. Its largest steering correction is 0.662 degrees.
Ordinary damping toward the anchor reduces error: half the correction has
2.611-mm disagreement. However it leaves terminal-cone residual 0.020728,
so it does not solve the constrained problem. The inherited slack budget
alone cannot certify that this shifted base is feasible.

\textbf{Fresh trajectory.} The Zhai-inspired potential guidance predicts the
moving target and produces an executable-model rollout, but its nearest node
at absolute $t=1.075$ s has 17.536-cm gap. Its terminal cone is feasible.
The maximum constraint violation is consequently the 2.464-cm gap shortage in
a hard completion-stage collision row. Guidance uses its own 0.7-m influence
padding; it does not read \texttt{collision.safetyMarginMeters}. Increasing
the optimization threshold therefore does not directly widen the freshly
generated path at a fixed initial state. The influence padding is not a promised
rectangle clearance after tracking the desired heading with vehicle dynamics.

The PCBF primary solve finds zero prefix slack and a convex-feasible point,
but this point already has 17.439-mm nonlinear disagreement. The largest
steering correction is only 0.536 degrees. A small input correction over many
stages is not by itself a bound on accumulated position error.
The CLF step increases the maximum to 26.947 mm. The current line search is
\[
 z(\alpha)=z^{P}+\alpha(z^{C}-z^{P}),
\]
so $\alpha\to0$ approaches the \emph{PCBF point}, not the original nonlinear
anchor. Its three actual trials yield 19.701, 18.003 and 17.579 mm; the limiting
primary point remains above 10 mm. Conversely, halfway from the original
anchor to the primary point gives 4.288-mm disagreement but leaves a
1.121-cm affine-constraint deficit. Accuracy and feasibility must both hold.

\textbf{Error accumulation.} All three nontrivial maxima in the table occur
2.45 s into the prediction, at absolute $t=2.70$ s. At the near-encounter
node $t=1.10$ s, shifted CLF, fresh PCBF and fresh CLF errors are respectively
3.086, 5.158 and 6.880 mm. The maximum is an accumulated future-model discrepancy,
not a centimeter-scale immediate actuation error. At the maximum, translation
contributes 9.520/15.887/24.295 mm and the heading term contributes
1.015/1.552/2.652 mm. This diagnostic does not justify dropping the future
accuracy requirement or relaxing a physical margin.

\textbf{Design implication.} Improving initialization remains useful: it should
respond to the requested margin and reduce the repair needed at the predicted
encounter while retaining its role as a search seed. It need not be a safety
certificate. The associated PCBF repair should also favor a small displacement
from that seed; minimizing safety slack alone does not select the smallest
nonlinear-model perturbation among equally feasible results. Such changes could
help provide a convex-feasible, model-accurate damping base, but no proposed
change is implemented or validated by this diagnostic. The measured failure
cannot be assigned entirely to a bad avoidance side, an infeasible physical
maneuver, or an overly large first steering jump.

\textbf{Artifacts.} Captured primary/secondary points and complete assembled
problems reside in the diagnostic cache. Compact results and source/input
provenance are in \path{report/INITIALIZATION_ERROR_DIAGNOSIS_20261003/}.
\end{document}
